\chapter{The Cost Monad}
\label{ch:costmonad}

\section{What is being built}

Given a continued interactive GSLT $G$, the cost endofunctor $\Cost$ produces a calculus
$\Cost(G)$ in which every interaction is gated on the consumption of a token. The design
question is where the gating lives, and the answer that makes the construction work is:
\emph{in the grammar of the image}.

In $\Cost(G)$ the functor adjoins a sort $\SigCat$ of signatures, a wrapped-term sort $\Tw$
with the wrapper $\{\cdot\}_s : P \times \SigCat \to \Tw$, and a sort $\Stk$ of token stacks
\[
  S ::= (\,) \mid s :: S, \qquad s \in \SigCat .
\]
The pivotal discipline is grammatical: in $\Cost(G)$ every continuation slot of $K_p$ and
$K_e$ has sort $\Tw$. A continuation in the metered calculus is therefore not a bare
process awaiting metering; it is a \emph{metered thunk} $\{P\}_s$, whose activation is
gated by a token keyed to $s$, and whose own continuations are again thunks.

The signature constructor is $\# = \mathrm{digest} \circ \cf$: the section of
Definition~\ref{ob:def:section} chooses a canonical representative, and a collision-resistant
digest commits to it. Note that $\#$ does \emph{not} respect $\equiv$, and must not,
since a commitment identifying congruent-but-distinct representatives would be no
commitment. Signatures never reduce; they are a rewrite-free sub-theory meeting the
behavioural theory only at the matching guard, so the $\equiv$-incoherence of $\#$ is
harmless for bisimulation, which factors as ``$P$-bisimulation gated by key matching.''

\section{The gated rules}

A wrapped redex is forced by consuming a matching token. The single-signature rule seals
the whole redex under one $s$ and funds it from a purse located at the surface
$L = \near(I,J)$:
\[
  \frac{L = \near(I,J) \qquad
        (\{C_p'\}_{t_1}, \{C_e'\}_{t_2}) = \compute(I,J,\{C_p\}_{t_1},\{C_e\}_{t_2})}
       {K\bigl(\{K(K_p(I,\{C_p\}_{t_1}), K_e(J,\{C_e\}_{t_2}))\}_s,\ S(L, s::p)\bigr)
        \longrightarrow
        K\bigl(K'(\{C_p'\}_{t_1}, \{C_e'\}_{t_2}),\ S(L,p)\bigr)}
  \tag{R1}
\]
The remaining cases seal the two interaction surfaces rather than the whole redex,
funding per surface, and are what the associative--commutative structure of $\mid$
requires when a redex is split across the bag or its tokens are supplied separately. We
write them as (R2) and (R3) and do not reproduce them here.

The decisive point is what is \emph{absent}. There is no re-wrapping step and no lifted
contraction. By the typing of $\compute$, the residual $K'(\{C_p'\}_{t_1},
\{C_e'\}_{t_2})$ is already built from wrapped terms.

\begin{definition}[Well-wrapped]
A configuration of $\Cost(G)$ is \emph{well-wrapped} when it is well-sorted in the
grammar of $\Cost(G)$ --- in particular, every continuation occupies a slot of sort
$\Tw$, so every redex occurrence lies inside a wrapper.
\end{definition}

\begin{lemma}[Subject reduction for wrapping]
Well-wrapping is preserved by (R1)--(R3).
\end{lemma}

\begin{proof}[Proof sketch]
Each rule replaces a wrapped redex by $K'(C_p',C_e')$ with $(C_p',C_e') = \compute(\dots)$.
Since $\compute$ lands in $\Tw \times \Tw$ and $K'$ packages terms of sort $\Tw$, the
residual is well-sorted; the consumed cell leaves a well-sorted stack; redexes elsewhere
are untouched.
\end{proof}

\begin{corollary}[No leak, by construction]
In a well-wrapped configuration no redex can fire without being unwrapped, i.e.\ without
consuming a token. There is no dynamic wrapping operation; the cost-accounting steps only
ever unwrap.
\end{corollary}

\begin{remark}[Eager versus lazy metering]
An alternative discipline re-wraps the contractum at contraction time, traversing it to
sign each exposed redex --- charging eagerly for every redex a step exposes. Wrapping by
construction charges lazily: each redex is born wrapped and is charged only when actually
forced. Lazy metering pays for work performed, not work merely exposed; redexes beneath a
discarded binder are never charged. The token stack drains exactly in step with the
reductions performed, which is also the operational reality of fuel consumption.
\end{remark}

\begin{remark}[Duplication needs no fresh signatures]
Because multiplicity is carried by the stack and not by key distinctness, copying a
wrapped term $\{Q\}_s$ copies its key but not its tokens: $n$ copies of an $s$-redex
require $n$ cells keyed to $s$. Created redexes --- where a substituted wrapped term
lands in head position --- are guarded for the same reason, since the substituted
material carries its own wrapper. One mechanism handles duplicated and created redexes
alike.
\end{remark}

\section{Two monoids: space and time}

The construction places two combination operators side by side, of different character.

The interaction constructor $K$ is \emph{spatial}: it records what is in contact with
what, and may carry equations --- associative--commutative in rho, free in $\lambda$.

The cons operator $::$ is \emph{temporal}: it records what is spent next, then next. It
is a free monoid, never commutative, because reduction is sequential even when
interaction is parallel. Each forcing pops a cell; the order of popping is the order of
steps. The stack is the time axis of the computation reified as a term: $\equiv$ leaves
it invariant, $\longrightarrow$ pops it, and its ordering is the arrow of the computation.

\begin{proposition}[Stack consumption is the modulus]
For a terminating reduction in $\Cost(G)$, the number of cells consumed equals the number
of forced redexes, which equals the number of cuts eliminated --- including those
introduced by duplication, each charged when forced. The consumed prefix of the temporal
stack is an exact operational modulus for the cut elimination, realised by running the
reduction rather than by a separate traversal.
\end{proposition}

\section{The monad and its resolution}

\begin{construction}[Unit and multiplication]
Define $\eta_G(P) = \{P\}_{(\,)}$, wrapping each term with the unit signature. Since
$(\,)$ is the unit of the signature monoid, a $(\,)$-keyed redex fires without net
resource, so $\eta_G$ embeds $G$ as the cost-free fragment of $\Cost(G)$. Define
$\mu_G : \Cost^2(G) \to \Cost(G)$ by multiplying nested signatures,
$\{\{P\}_{s_2}\}_{s_1} \mapsto \{P\}_{s_1 * s_2}$, and concatenating the stack-of-stacks
into a single stack.
\end{construction}

\begin{proposition}[The cost monad]
$(\Cost, \eta, \mu)$ is a monad on $\catciGSLT$. Its laws descend from the laws of the two
constituent monoids: the unit laws from $s * (\,) = s = (\,) * s$ together with the
singleton/empty laws of the list monad, and associativity from associativity of $*$ and
of concatenation.
\end{proposition}

\begin{remark}[A genuine, non-idempotent resource monad]
$\Cost$ is not idempotent: $\mu_G$ is no isomorphism, because concatenating two stacks
forgets the boundary between them and multiplying two grades forgets the factorisation.
Flattening is a true merge of two resource accounts, not the collapse of a redundant
copy. This is the expected signature of a resource monad, and it distinguishes $\Cost$
from idempotent closure monads.
\end{remark}

\begin{proposition}[Free--forgetful resolution]
\label{ob:prop:cost-adjunction}
Let $\mathrm{Cost}\text{-}\catciGSLT$ be the category of cost-accounted continued
interactive GSLTs. The functor $\mathrm{Install}$ that adjoins the apparatus is left
adjoint to the functor $\mathrm{Forget}$ that strips it, and the induced monad on
$\catciGSLT$ is $\Cost$.
\end{proposition}

The embedding is structural and strict: $\mathrm{Install}$ adjoins sorts and rules,
$\mathrm{Forget}$ erases them on the nose. Crucially, forgetting is not
behaviour-preserving in the metering sense --- erasing the apparatus removes the gating,
so the forgotten theory runs its interactions without friction. The resolution answers
``what apparatus was added?'' and lets us remove it; it does not claim the metered and
un-metered dynamics agree.

\begin{remark}[Internalisation]
\label{ob:rem:internalise-cost}
When the base theory is sufficiently expressive --- enough to code its own higher-order
syntax, whether by binding or by reflection --- the entire metering apparatus is itself
expressible in the base. Tokens become encoded data and the gated rules become an
interpreter loop, so that the base performs the metering with its own computation. This
is a second and quite different erasure from forgetting: rather than stripping the
apparatus it dissolves it into the base. We record its availability and do not develop
it here.
\end{remark}

\section{Capabilities and located purses}
\label{ob:subsec:capabilities}

A cost-accounted calculus is a calculus in which the right to draw on a resource stack is
itself an authority. We must therefore ask what confers that authority, and ensure the
construction does not silently grant it. The answer turns entirely on the equational
theory of $K$, in the manner already anticipated in Section~\ref{ob:subsec:surfaces}.

\paragraph{Free $K$: position is the capability.} When $K$ carries no equations, nothing
can drift into contact with a stack it does not already neighbour. A capability to draw
on a stack $S$ is then literally the context $K([\,], S)$ --- the right to occupy the hole
adjacent to $S$. Because $K$ is rigid, occupying that position is the whole of the
authority, and confinement is automatic. This is the $\lambda$-calculus regime and it is
the object-capability ideal: authority is structural, held by reference, and unforgeable
by construction.

\paragraph{Non-free $K$: the leak.} The moment $K$ carries any equation, position is no
longer exclusive. An equation is a congruence step that changes adjacency without firing
a cut, so it lets a program drift into contact with a stack it did not structurally
neighbour. Associative--commutativity is the limit: the bag is a freely mixed soup, every
program adjacent to every stack, and rho lives exactly here. Taking proximity as
authority in that regime is precisely \emph{ambient authority}~\cite{MillerYeeShapiro2003},
and it is a leak.

\paragraph{The repair: nominal surfaces and located stacks.} To recover an
object-capability discipline under a non-free $K$, carry the surface nominally and gate on
a nearness operator $\near(I,J)$, which when defined returns the surface at which the two
meet. Interaction is licensed not by position but by the matchability of the surfaces.
Then locate the stack at a surface,
\[
  S(I,\ s :: S') ,
\]
so that the cookie jar carries a label naming which hands may reach in. Authority over a
resource is thereby named in the term rather than conferred by the geometry of the bag.

\begin{remark}[Nearness bounds but does not eliminate competition]
$\near(I,J)$ does not by itself make access exclusive: two receivers may both be near the
same surface. What it guarantees is that only holders of a matching surface may compete
at all --- it replaces ``anyone adjacent'' with ``anyone holding a matching surface.''
Exclusivity, where required, is a matter for the type discipline of
Chapter~\ref{ch:hypercube}.
\end{remark}
